% ================================================================================================ cómo leer
\chapter*{Antes de empezar}
\markboth{Antes de empezar}{}
\addcontentsline{toc}{chapter}{Antes de empezar}

TurbuLENS nació en la asignatura Técnicas Observacionales de la Maestría en Astronomía del Observatorio Astronómico
Nacional (OAN). Queríamos medir qué tan turbulento es el cielo de Bogotá con un telescopio de aficionado, una cámara de
fotos y una máscara de cartón con dos orificios, y nos encontramos con que los programas que había daban números que no
se podían defender. Así que escribimos uno, lo probamos con videos de valores conocidos y lo usamos en el OAN. Este
manual cuenta cómo usarlo y por qué funciona.

\section*{Cómo está organizado}

\begin{center}\small
\begin{tabularx}{\linewidth}{@{}>{\bfseries\color{acento}}l Y@{}}
Parte I   & \textbf{Primeros pasos.} Descarga, primera apertura, la ventana y una medición completa en diez minutos.
            Si nunca has abierto TurbuLENS, empieza aquí.\\
Parte II  & \textbf{Usar TurbuLENS.} Cada modo, pantalla por pantalla, con capturas numeradas. Incluye un capítulo
            entero para quien mide con su propio telescopio y su propia cámara (no hace falta ser del OAN).\\
Parte III & \textbf{Fundamentos.} La teoría justa para entender los resultados: turbulencia, el DIMM, la escala de placa,
            estadística, calibración de imágenes y el método satelital.\\
Parte IV  & \textbf{Por dentro.} Los fragmentos del código que hacen las mediciones, con sus números de línea reales,
            para quien quiera revisar, citar o mejorar el método.\\
Apéndices & Equipos del OAN, historial de versiones, créditos, licencia, cómo citar y referencias.
\end{tabularx}
\end{center}

\section*{Convenciones}

\begin{center}\small
\begin{tabularx}{\linewidth}{@{}l Y@{}}
\toprule
Se ve así & Significa\\
\midrule
\botonp{Generar}, \boton{Cancelar} & Un botón de la aplicación (los verdes son la acción principal de la pantalla).\\
\tecla{Mayús}\,+\,\tecla{→} & Teclas que se pulsan juntas.\\
\menu{Recorte y tramo} & Una página del menú de la izquierda o una opción de una lista.\\
\campo{Seguir a la estrella} & Una casilla, un campo o el título de una tarjeta.\\
\ruta{evaluadas\_dia\_…/} & Un archivo o una carpeta.\\
\marcanum{3} & Una marca numerada en una captura; la leyenda debajo de la figura dice qué es.\\
\bottomrule
\end{tabularx}
\end{center}

\begin{nota}[Los botones de ayuda.]
Casi todas las tarjetas de TurbuLENS tienen un botón redondo con un signo de exclamación. Al pulsarlo aparece una
explicación corta de esa parte de la pantalla. Este manual cuenta lo mismo con más calma y con ejemplos; los dos se
complementan.
\end{nota}

Los números con decimales se escriben con coma (1,242″/px), como en la aplicación en español. Los segundos de arco se
escriben con el símbolo ″: $1'' = 1/3600$ de grado, el tamaño aparente de una moneda de un centímetro vista a dos
kilómetros.
